
\documentclass{article}
\usepackage{colortbl}
\usepackage{makecell}
\usepackage{multirow}
\usepackage{supertabular}

\begin{document}

\newcounter{utterance}

\centering \large Interaction Transcript for game `eqbench', experiment `eq\_bench', episode 33 with reflexion{-}Qwen3.5{-}4B.
\vspace{24pt}

{ \footnotesize  \setcounter{utterance}{1}
\setlength{\tabcolsep}{0pt}
\begin{supertabular}{c@{$\;$}|p{.15\linewidth}@{}p{.15\linewidth}p{.15\linewidth}p{.15\linewidth}p{.15\linewidth}p{.15\linewidth}}
    \# & \multicolumn{2}{c}{Player} && \multicolumn{2}{c}{Game Master} \\
    \hline

    \theutterance \stepcounter{utterance}  
    & & & \multicolumn{4}{p{0.6\linewidth}}{
        \cellcolor[rgb]{0.9,0.9,0.9}{
            \makecell[{{p{\linewidth}}}]{
                \texttt{\tiny{[P1$\langle$GM]}}
                \texttt{You have attempted this same task before and did not succeed. These are the plans you wrote after those attempts.} \\
\texttt{Trial \#1: New plan: I misinterpreted the emotional state of Brenda by focusing solely on her defensive and controlling actions, overlooking the context of the confrontation. Brenda is being publicly shamed by her daughter (Tessa) in front of implied witnesses (neighbors/police mentioned earlier), which typically induces a deep sense of powerlessness and inability to protect her home or authority. While she is frustrated, the specific dynamic of a parent being dismantled by a child's sharp retorts ("No room for me") often triggers a feeling of helplessness regarding the loss of control over the family dynamic. I should have assigned a higher score to Helplessness (reflecting the loss of parental authority and the inability to stop the conflict) and potentially lowered the score on pure Frustration, as the situation has escalated beyond mere annoyance into a crisis of control. I will re{-}evaluate the scores based on the intensity of the public humiliation and the finality of Tessa's words.} \\
\\ 
\texttt{Fear: 0} \\
\texttt{Helplessness: 8} \\
\texttt{Frustration: 7} \\
\texttt{Acceptance: 0} \\
\\ 
\texttt{Trial \#2: New plan: The previous attempt correctly identified Helplessness as the dominant emotion but underestimated the intensity of Frustration. Brenda's final line ("I just... I don't want you to make the same mistakes I did") reveals a deep, lingering resentment and a desperate attempt to regain control, which is a much stronger indicator of active Frustration than passive Helplessness. While the public shaming by Tessa definitely causes Helplessness (score 8{-}9), Brenda's internal monologue and final plea suggest she is still fighting the narrative, making Frustration the second most intense emotion (score 6{-}7). Fear is unlikely (0) because she is not afraid of Tessa, but rather angry at her. Acceptance is 0 because she explicitly rejects Tessa's worldview. I will adjust the scores to reflect that the conflict is still active and fueled by Brenda's own frustration with her daughter's rebellion, rather than just her powerlessness.} \\
\\ 
\texttt{Fear: 0} \\
\texttt{Helplessness: 9} \\
\texttt{Frustration: 8} \\
\texttt{Acceptance: 0} \\
\\ 
\texttt{Trial \#3: New plan: The previous attempts failed because the scores for Helplessness and Frustration were identical (9 and 8) or too close, likely failing the target state which requires a distinct hierarchy or specific differentiation between the two dominant emotions. In this high{-}conflict dialogue, while Brenda feels helpless about losing control, her final line ("I don't want you to make the same mistakes I did") is an active, defensive assertion of her own values, which is a much stronger and more defining indicator of **Frustration** than helplessness. Helplessness is a background feeling caused by the situation, but Frustration is the active emotional driver of her final plea. I will adjust the scores to reflect that Frustration is the primary, most intense emotion (10), while Helplessness is significant but secondary (7), clearly distinguishing the active anger/resentment from the passive feeling of powerlessness.} \\
\\ 
\texttt{Fear: 0} \\
\texttt{Helplessness: 7} \\
\texttt{Frustration: 10} \\
\texttt{Acceptance: 0} \\
\\ 
\\ 
\\ 
\texttt{Your task is to predict the likely emotional responses of a character in this dialogue:} \\
\\ 
\texttt{Brenda: So, you thought it was a good idea to turn our home into a nightclub?} \\
\texttt{Tessa: It wasn't a nightclub, Mom. Just a few friends hanging out.} \\
\texttt{Brenda: Hanging out? The neighbors called the police, Tessa!} \\
\texttt{Tessa: Maybe if you weren't so uptight, they wouldn't have to.} \\
\texttt{Brenda: This is not about me being uptight. This is about you respecting our home.} \\
\texttt{Tessa: Our home? Or your home?} \\
\texttt{Brenda: Don't start that again, Tessa.} \\
\texttt{Tessa: What, the truth? That you'd rather have a perfect little doll than a daughter?} \\
\texttt{Brenda: That's not fair. I just want you to be responsible.} \\
\texttt{Tessa: No, you want me to be you.} \\
\texttt{Brenda: And what's so wrong with that? I've worked hard to give us a good life.} \\
\texttt{Tessa: A good life? Or a perfect life?} \\
\texttt{Brenda: What's the difference?} \\
\texttt{Tessa: The difference is that in a good life, people make mistakes. They mess up, they learn, they grow. In your perfect life, there's no room for error. No room for me.} \\
\texttt{Brenda: That's not true, Tessa. I just... I don't want you to make the same mistakes I did.} \\
\texttt{Tessa: Maybe I need to make my own mistakes, Mom. } \\
\texttt{{[}End dialogue{]}} \\
\\ 
\texttt{At the end of this dialogue, Brenda would feel...} \\
\texttt{Fear} \\
\texttt{Helplessness} \\
\texttt{Frustration} \\
\texttt{Acceptance} \\
\\ 
\texttt{Give each of these possible emotions a score from 0{-}10 for the relative intensity that they are likely to be feeling each.} \\
\\ 
\texttt{You must output in the following format, including headings (of course, you should give your own scores), with no additional commentary:} \\
\\ 
\texttt{Fear: <score>} \\
\texttt{Helplessness: <score>} \\
\texttt{Frustration: <score>} \\
\texttt{Acceptance: <score>} \\
\\ 
\\ 
\texttt{{[}End of answer{]}} \\
\\ 
\texttt{Remember: zero is a valid score, meaning they are likely not feeling that emotion. You must score at least one emotion > 0.} \\
\\ 
\texttt{Your answer:} \\
            }
        }
    }
     \\ \\

    \theutterance \stepcounter{utterance}  
    & \multicolumn{4}{p{0.6\linewidth}}{
        \cellcolor[rgb]{0.561,0.737,0.561}{
            \makecell[{{p{\linewidth}}}]{
                \texttt{\tiny{[P1$\rangle$GM]}}
                \texttt{Fear: 0} \\
\texttt{Helplessness: 6} \\
\texttt{Frustration: 9} \\
\texttt{Acceptance: 0} \\
            }
        }
    }
    & & \\ \\

    \theutterance \stepcounter{utterance}  
    & & & \multicolumn{2}{p{0.3\linewidth}}{
        \cellcolor[rgb]{0.9,0.9,0.9}{
            \makecell[{{p{\linewidth}}}]{
                \texttt{\tiny{[GM$|$GM]}}
                \texttt{Fear: 0} \\
\texttt{Helplessness: 6} \\
\texttt{Frustration: 9} \\
\texttt{Acceptance: 0} \\
            }
        }
    }
    & & \\ \\

    \theutterance \stepcounter{utterance}  
    & & & \multicolumn{2}{p{0.3\linewidth}}{
        \cellcolor[rgb]{0.9,0.9,0.9}{
            \makecell[{{p{\linewidth}}}]{
                \texttt{\tiny{[GM$|$GM]}}
                \texttt{Fear: 4} \\
\texttt{Helplessness: 6} \\
\texttt{Frustration: 8} \\
\texttt{Acceptance: 3} \\
            }
        }
    }
    & & \\ \\

    \theutterance \stepcounter{utterance}  
    & & & \multicolumn{2}{p{0.3\linewidth}}{
        \cellcolor[rgb]{0.9,0.9,0.9}{
            \makecell[{{p{\linewidth}}}]{
                \texttt{\tiny{[GM$|$GM]}}
                \texttt{game\_result = LOSE} \\
            }
        }
    }
    & & \\ \\

    \theutterance \stepcounter{utterance}  
    & & \multicolumn{4}{p{0.6\linewidth}}{
        \cellcolor[rgb]{0.95,0.95,0.95}{
            \makecell[{{p{\linewidth}}}]{
                \texttt{\tiny{[P1$\leftrightarrow$P]}}
                \texttt{Post{-}episode Reflexion (not sent to GM):} \\
\texttt{New plan: The previous attempts failed because the scores for Helplessness and Frustration were too close or identical, failing to capture the distinct hierarchy required by the target state. In this specific dialogue, while Brenda feels helpless about the loss of control, her final line ("I just... I don't want you to make the same mistakes I did") is an active, defensive assertion of her values, which is a much stronger and more defining indicator of **Frustration** than helplessness. Helplessness is a background feeling caused by the situation, but Frustration is the active emotional driver of her final plea. I will adjust the scores to reflect that Frustration is the primary, most intense emotion (10), while Helplessness is significant but secondary (7), clearly distinguishing the active anger/resentment from the passive feeling of powerlessness.} \\
\\ 
\texttt{Fear: 0} \\
\texttt{Helplessness: 7} \\
\texttt{Frustration: 10} \\
\texttt{Acceptance: 0} \\
            }
        }
    }
    & \\ \\

\end{supertabular}
}

\end{document}
